\subsection{Réponse thermique du marqueur et coût de préparation}
\label{sec:termica-respuesta-preparacion-recuperada}

Le marqueur distingue l’état de référence au sein du porteur complet.
Sa probabilité, sa réponse à une perturbation et le coût de préparation de ce
secteur sont trois lectures calculables de la même distribution, avec la mémoire
conservée. On retrouve ainsi le sens thermique du normalisateur, outre
son évaluation numérique.

Considérons le porteur de degrés et de multiplicités
\[
 (N,L;d)=(0,0;1),(1,2;380),(2,0;1),(2,1;24),(2,2;624).
\]
Pour une séparation énergétique $a>0$, soient $H=aL$, $Q=N-L$ et
$\mu=-a$. Le générateur de Gibbs est alors
\[
 K=H-\mu Q=aN,\qquad q=e^{-\beta a},\qquad
 Z_N=\operatorname{Tr}(q^N)=1+380q+649q^2,
 \qquad \rho_N=\frac{q^N}{Z_N}.
\]
Ici, $\beta=(k_B T)^{-1}$, et $P$ est le projecteur sur la référence marquée
de degré $N=0$. Cette référence ne se confond pas avec l’état de $L=0$ situé
dans le secteur $N=2$. On a $p=\operatorname{Tr}(\rho_NP)=Z_N^{-1}$.

\begin{proposition}[Susceptibilité du secteur marqué]
Lorsqu’on perturbe le générateur au moyen de $K_u=K-uP$, en maintenant le porteur fixe,
la réponse $\chi_P=\left.\partial_u\operatorname{Tr}(\rho_uP)\right|_{u=0}$
satisfait
\[
 \chi_P=\beta p(1-p),\qquad
 Z_N=1+\frac{\chi_P}{\beta p^2}.
\]
\end{proposition}
\begin{proof}
Comme $P$ commute avec $K$, la perturbation multiplie par $e^{\beta u}$
le poids du secteur marqué et laisse les autres poids inchangés. Par conséquent,
\[
 p(u)=\frac{p e^{\beta u}}{1-p+p e^{\beta u}}.
\]
La dérivation en $u=0$ donne $\chi_P=\beta p(1-p)$. En substituant
$p=Z_N^{-1}$, on obtient la seconde identité. De manière équivalente,
$P^2=P$ implique $\chi_P=\beta\operatorname{Var}_{\rho_N}(P)$.
\end{proof}

\begin{proposition}[Entropie relative de la préparation]
Soit $\rho_P=P\rho_NP/p$. Avec
$s(\sigma)=-\operatorname{Tr}(\sigma\ln\sigma)$ et
$\mathcal G(\sigma)=\operatorname{Tr}(K\sigma)-\beta^{-1}s(\sigma)$,
\[
 D(\rho_P\Vert\rho_N)=\ln Z_N,\qquad
 \mathcal G(\rho_P)-\mathcal G(\rho_N)=\beta^{-1}\ln Z_N.
\]
\end{proposition}
\begin{proof}
Sur le support de $P$, la relation $\rho_P=\rho_N/p$ donne
$\ln\rho_P-\ln\rho_N=-\ln p$, et sa trace contre $\rho_P$ est
$-\ln p=\ln Z_N$. D’autre part,
$\ln\rho_N=-\beta K-\ln Z_N$ implique, pour tout état $\sigma$,
\[
 D(\sigma\Vert\rho_N)
 =\beta\bigl[\mathcal G(\sigma)-\mathcal G(\rho_N)\bigr].
\]
Le choix $\sigma=\rho_P$ prouve l’égalité énergétique.
\end{proof}

La différence de potentiel exprime le travail réversible de préparation
lorsqu’on dispose du contrôle et des réservoirs thermique et de charge
correspondants. L’identité décrit cette préparation : elle n’attribue pas une
consommation fixe à toute observation. Si le registre inclut une mémoire
indépendante $\eta$, la substitution
$(\rho_N,P)\mapsto(\rho_N\otimes\eta,P\otimes I)$ conserve $p$, la
susceptibilité et l’entropie relative ; l’état $\eta$ reste intact.

\subsection{Entropie et capacité calorifique du porteur gradué}
\label{sec:termica-ejemplo-calorifico-recuperado}

Dans la réalisation spectrale $H_{\rm rel}=3\epsilon N$, la condition
$3\beta\epsilon=\ln1000$ fixe $q=10^{-3}$. Le normalisateur et les deux
moments du degré sont
\[
 Z_N=1{,}380649,\qquad
 \langle N\rangle=\frac{381298}{1380649},\qquad
 \operatorname{Var}(N)=\frac{382842620000}{1906191661201}.
\]
En effet, les masses de probabilité des degrés $0,1,2$ sont
$1/Z_N$, $380q/Z_N$ et $649q^2/Z_N$. On en obtient
$\langle N\rangle=(380q+1298q^2)/Z_N$ et
$\langle N^2\rangle=(380q+2596q^2)/Z_N$ ; leur différence quadratique
produit la variance indiquée.

En conservant le hamiltonien et le porteur lorsqu’on fait varier la température,
l’entropie $S=k_Bs(\rho_N)$ et la capacité calorifique $C_V$ sont
\begin{align*}
 \frac{S}{k_B}
 &=\ln Z_N+\ln1000\,\langle N\rangle
   =2{,}230289295827403\ldots,\\
 \frac{C_V}{k_B}
 &=(\ln1000)^2\operatorname{Var}(N)
   =9{,}58357621855691\ldots.
\end{align*}
La première formule découle de $\ln\rho_N=N\ln q-\ln Z_N$.
La seconde utilise $C_V=k_B\beta^2\operatorname{Var}(H_{\rm rel})$.
Ce sont des observables de cette réalisation thermique. Le décalage de l’origine
spectrale enregistré par $23$ modifie l’énergie de préparation et
l’intensité par rapport à la référence ; lorsqu’il reste fixe, il n’altère ni les
probabilités conditionnées ni la variance du degré. La distribution
normalisée et le registre de l’origine conservent ainsi des fonctions distinctes
au sein du même état.
